\begin{abstract}
\noindent
Physics-informed neural networks (PINNs) and neural--mechanistic hybrid models
estimate the parameters of signalling models, and their failures are
attributed either to the estimator or to non-identifiability.
They call for opposite remedies and are rarely told apart. The Cram\'er--Rao bound of the experiment performed separates
them. We carry this accounting through a
seven-state model of WNT, retinoic acid and HOX signalling in colorectal cancer,
from healthy crypt to severe APC loss. Three inverse PINN formulations recover 2--21 of 34
parameters per regime, but the bound says an efficient estimator would recover
132.5 of 136 regime--parameter pairs, a converged maximum-likelihood fit
attains it on independent datasets where the data are informative, and the PINNs' errors are 10--92
times the bound at the median: an estimator gap. A Schur-complement
Gauss--Newton PINN, which freezes the network's features so that the problem
becomes separable with a block-arrowhead structure, closes the gap wherever its
continuation converges (both low-WNT regimes on two datasets), there matching
the classical estimator and reporting its own precision. Where the bound itself
is wide no estimator helps: APC functionality, the parameter that defines the
disease stage, loses three orders of magnitude of precision (the identifiable
product it enters, thirty-fold) as \bcat{} saturates, and one WNT knockdown restores it while information-matched
conditions do not. A learned regulatory term opens an analogous gap: its
constraint $f(0)=0$ is informative only if a condition drives the regulator to
zero, so the host equation's basal rate is lost under every network
parameterisation tried. A Fisher analysis of the hybrid equation, a
fixed-effort dose-response and a prospective test of a training-free design
table show that a knockdown visiting the anchor, in a condition that still
contains the basal term, removes most of the loss. The first diagnostic should be the design's information budget.
\end{abstract}

\section*{Author summary}
Mathematical models of cell signalling are only as useful as the numbers that go
into them, and those numbers, rates of production and degradation, are usually
estimated from experiments. Machine-learning methods that combine neural
networks with the model's equations promise to do this from sparse data, but
when an estimate comes out wrong it is hard to know why: was the method too
weak, or did the experiment simply not contain the information? We show that a
classical quantity from statistics, the Cram\'er--Rao bound, answers this
question before any network is trained, and we use it to work through a model of
the signalling network that goes awry in colorectal cancer. For most parameters
the information was there and a better solver, which we develop, recovers them.
For the parameter that marks the stage of the disease, and for parameters that
sit next to a term learned by a network, the information was missing, and only a
different experiment, a targeted knockdown, brought it back. The bound tells an
experimentalist which of the two to change.
